\section*{\fontsize{16pt}{15pt}\selectfont DESIGN AND IMPLEMENTATION OF ACL \\IN 3-LAYER NETWORK MODEL WITH DMZ \\ABSTRACT}

\begin{itemize}
    \item[-] \textbf{Research issues:}
    This report focuses on researching and implementing Access Control List (ACL) techniques within a hierarchical 3-layer network model (Core - Distribution - Access) integrated with a DMZ. The main content includes:
        \begin{itemize}
        \item[•] Chapter 1: Design and Implementation of ACL in a 3-layer network model with DMZ.
        \end{itemize}
   
    \item[-] \textbf{Approach:}
        \begin{itemize}
        \item[•] \textbf{Theory:} Understanding the principles of Standard ACL and Extended ACL, packet filtering rules based on IP addresses, ports, and protocols.
        \item[•] \textbf{Practice:} Using Mininet to simulate the network system and iptables tools to configure practical security policies.
    \end{itemize}
    
    \item[-] \textbf{Solution methodology:}
    Based on the security requirements of a hypothetical university network (TDTU), the team analyzes traffic flows to design corresponding allow/deny rules. Python scripts are used to automate topology construction and configuration application.

    \item[-] \textbf{Key results achieved:}
    Successfully built a network model with full segmentation: Internal (Inside), Public (DMZ), and Internet (Outside). Ensured security policies: preventing unauthorized access from the Internet to the internal network, isolating student and lecturer networks, and protecting device management rights.
\end{itemize}

\newpage